\documentclass[../../main.tex]{subfiles}
\begin{document}
% SOURCE p=212; printed=197; object=ch06-distribution-title; region=middle
\section{分布函数与弱 $L^p$}\label{section:ch06-distribution-weak-lp}
% SOURCE p=212; printed=197; object=ch06-distribution-definition; region=middle
\begin{definition}[分布函数]\label{definition:ch06-distribution-function}
若 $f$ 为测度空间 $(X,\mathscr M,\mu)$ 上的可测函数，定义其\textbf{分布函数} $\lambda_f:(0,\infty)\to[0,\infty]$ 为
\begin{align*}
\lambda_f(\alpha)=\mu\{x:|f(x)|>\alpha\}.
\end{align*}
\end{definition}
% SOURCE p=212; printed=197; object=ch06-distribution-name-comparison; region=middle
它与\cref{section:ch01-borel-measures-real-line,section:ch10-probability-basic}中讨论的“分布函数”密切相关，但并不相同。下面用一个命题汇集其基本性质。
% SUPPLEMENT: id=v14-ch06-distribution-motivation; parent=ch06-distribution-definition; kind=motivation+hypotheses; origin=independent-derivation
\begin{remark}
\textbf{动机与逐条件说明（重构新增）。}分布函数记录函数超过各个高度的区域有多大，舍去了这些区域的位置，却保留了许多积分估计所需的信息。$f$ 可测保证水平集可测；取绝对值统一处理实值与复值函数；严格不等号与后面的右连续性相配。参数限于 $\alpha>0$，避免零水平直接包含全部非零支集；值域允许 $\infty$，因为这里不假设有限测度或 $f\in L^p$。等价代表的水平集至多相差零集，因此 $\lambda_f$ 只依赖几乎处处类。这里记录的是模值的上尾大小，通常的概率累积分布函数记录有序变量不超过某阈值的概率，二者不可混淆。
\end{remark}
% END SUPPLEMENT: id=v14-ch06-distribution-motivation
% SOURCE p=212; printed=197; object=ch06-distribution-properties; number=6.22; region=middle-lower
\begin{proposition}\label{proposition:ch06-distribution-properties}
\begin{enumerate}
\item\label{item:ch06-distribution-right-continuous} $\lambda_f$ 递减且右连续。
\item\label{item:ch06-distribution-order} 若 $|f|\leqslant|g|$，则 $\lambda_f\leqslant\lambda_g$。
\item\label{item:ch06-distribution-monotone-limit} 若 $|f_n|\uparrow|f|$，则 $\lambda_{f_n}\uparrow\lambda_f$。
\item\label{item:ch06-distribution-sum} 若 $f=g+h$，则 $\lambda_f(\alpha)\leqslant\lambda_g(\alpha/2)+\lambda_h(\alpha/2)$。
\end{enumerate}
\end{proposition}
% SOURCE p=212; printed=197; object=ch06-distribution-properties-proof; region=lower
\begin{proof}
令 $E(\alpha,f)=\{x:|f(x)|>\alpha\}$。当 $\alpha<\beta$ 时，$E(\alpha,f)\supset E(\beta,f)$，所以 $\lambda_f$ 递减。又 $E(\alpha,f)$ 为 $\{E(\alpha+1/n,f)\}_{n\geqslant1}$ 的递增并，故 $\lambda_f$ 右连续。若 $|f|\leqslant|g|$，则 $E(\alpha,f)\subset E(\alpha,g)$，因而 $\lambda_f\leqslant\lambda_g$。若 $|f_n|\uparrow|f|$，则 $E(\alpha,f)$ 为 $\{E(\alpha,f_n)\}$ 的递增并，故 $\lambda_{f_n}\uparrow\lambda_f$。最后，若 $f=g+h$，则
\begin{align*}
E(\alpha,f)\subset E(\alpha/2,g)\cup E(\alpha/2,h),
\end{align*}
从而 $\lambda_f(\alpha)\leqslant\lambda_g(\alpha/2)+\lambda_h(\alpha/2)$。
\end{proof}
% SUPPLEMENT: id=v14-ch06-distribution-properties-details; parent=ch06-distribution-properties; kind=proof-detail+motivation+hypotheses; origin=independent-derivation
\begin{remark}
\textbf{集合极限与条件说明（重构新增）。}“递减”在此指非增。右连续性使用测度对递增并的连续性，即使结果为无穷也成立；不需要用须有有限初始测度的递减交连续性。若 $|f(x)|>\alpha$，可选 $n$ 使 $1/n<|f(x)|-\alpha$，证明了原式中递增并确等于所需水平集；结合单调性，沿 $\alpha+1/n$ 的极限即为完整右极限。(c) 同理：严格超过 $\alpha$ 的极限值最终被某个 $|f_n|$ 严格超过。模序或单调收敛只需几乎处处成立，删去相应可数个例外零集即可。对于 (d)，若两项模都不超过 $\alpha/2$，则和的模不超过 $\alpha$；更一般地，对 $a,b>0$ 有 $\lambda_{g+h}(a+b)\leqslant\lambda_g(a)+\lambda_h(b)$。这些集合包含关系是以后估计弱空间和重排的基本工具；没有有限性或 $\sigma$ 有限性要求。
\end{remark}
% END SUPPLEMENT: id=v14-ch06-distribution-properties-details
% SOURCE p=212; printed=197; object=ch06-distribution-stieltjes; region=bottom
假设对每个 $\alpha>0$ 都有 $\lambda_f(\alpha)<\infty$。由\subcref{proposition:ch06-distribution-properties}{item:ch06-distribution-right-continuous}，$\lambda_f$ 确定 $(0,\infty)$ 上的一个负 Borel 测度 $\nu$，使得当 $0<a<b$ 时，
\begin{align*}
\nu((a,b])=\lambda_f(b)-\lambda_f(a).
\end{align*}
（\cref{section:ch01-borel-measures-real-line}在 $\mathbb R$ 上构造 Borel 测度的方法，同样适用于 $(0,\infty)$。）因此可以考虑 $(0,\infty)$ 上函数 $\phi$ 的 Lebesgue–Stieltjes 积分 $\int\phi\,\mathrm d\lambda_f=\int\phi\,\mathrm d\nu$。下一个结果说明，$X$ 上 $|f|$ 的函数的积分可化为这种 Lebesgue–Stieltjes 积分。
% SOURCE p=212; printed=197; object=ch06-distribution-pushforward; number=6.23; region=bottom
\begin{proposition}\label{proposition:ch06-distribution-pushforward}
若对所有 $\alpha>0$，$\lambda_f(\alpha)<\infty$，且 $\phi$ 为 $(0,\infty)$ 上的非负 Borel 可测函数，则
\begin{align*}
\int_X\phi\circ|f|\,\mathrm d\mu=-\int_0^\infty\phi(\alpha)\,\mathrm d\lambda_f(\alpha).
\end{align*}
\end{proposition}
% SOURCE p=213; printed=198; object=ch06-distribution-pushforward-proof; region=top
\begin{proof}
若 $\nu$ 为 $\lambda_f$ 确定的负测度，则
\begin{align*}
\nu((a,b])&=\lambda_f(b)-\lambda_f(a)
=-\mu\{x:a<|f(x)|\leqslant b\}\\
&=-\mu\bigl(|f|^{-1}((a,b])\bigr).
\end{align*}
由延拓的唯一性（\cref{theorem:ch01-measure-extension}），对每个 Borel 集 $E\subset(0,\infty)$ 都有 $\nu(E)=-\mu(|f|^{-1}(E))$。这正说明当 $\phi$ 为 Borel 集的示性函数时，所求公式成立，因而对简单函数也成立。一般情形由\cref{theorem:ch02-simple-approximation}与单调收敛定理得到。
\end{proof}
% SUPPLEMENT: id=v14-ch06-distribution-pushforward-details; parent=ch06-distribution-pushforward; kind=proof-detail+motivation+hypotheses; origin=independent-derivation
\begin{remark}
\textbf{零点约定、唯一性与条件说明（重构新增）。}原陈述把 $\phi$ 定义在 $(0,\infty)$，所以在 $f=0$ 处写 $\phi\circ|f|$ 需要补充约定：本式将 $\phi$ 以 $\phi(0)=0$ 延拓。等价地，左边只在 $\{|f|>0\}$ 上积分。原证明从 $E\subset(0,\infty)$ 的示性函数开始，正对应这一零延拓；若任意指定非零的 $\phi(0)$，公式一般不成立，例如非零有限测度空间上的 $f=0$。

令 $\eta(E)=\mu(|f|^{-1}(E))$，$E\subset(0,\infty)$。这是正 Borel 测度，而 $\rho=-\nu$ 也是正测度。它们在所有半开区间 $(a,b]$ 上相同，且在 $(1/n,n]$ 上均有限，因为该集合的逆像包含于 $\{|f|>1/n\}$。这些区间可数地覆盖正半轴，故所用唯一延拓确有 $\sigma$ 有限性保证；无须原空间 $X$ 本身 $\sigma$ 有限。于是 $\eta=\rho$。对非负简单 $\phi_n\uparrow\phi$，两边分别以正测度 $\mu,\rho$ 使用单调收敛，即得原公式，包括两边为 $+\infty$ 的情形。记号右边的负号应理解为对正测度 $-\nu$ 积分，避免任何无穷减无穷。

动机是把积分从原空间转移到模值所在的一维正半轴；水平集有限让 Stieltjes 增量不含 $\infty-\infty$，并提供唯一性所需的局部有限性。$\phi$ 的 Borel 可测性保证复合可测，非负性使无穷积分与单调逼近合法。以后若处理有符号 $\phi$，必须另加绝对可积等条件，不能直接照搬非负版本。
\end{remark}
% END SUPPLEMENT: id=v14-ch06-distribution-pushforward-details
% SOURCE p=213; printed=198; object=ch06-layer-cake-introduction; region=upper-middle
我们最关心的是 $\phi(\alpha)=\alpha^p$ 的情形，它给出
\begin{align*}
\int|f|^p\,\mathrm d\mu=-\int_0^\infty\alpha^p\,\mathrm d\lambda_f(\alpha).
\end{align*}
对右边分部积分（\cref{theorem:ch03-stieltjes-parts}），可得到更有用的形式 $\int|f|^p\,\mathrm d\mu=p\int_0^\infty\alpha^{p-1}\lambda_f(\alpha)\,\mathrm d\alpha$。除非已知 $\alpha^p\lambda_f(\alpha)$ 在 $\alpha\to0$ 与 $\alpha\to\infty$ 时均趋零，否则这一计算的合法性并不明显；不过结论确实正确。
% SOURCE p=213; printed=198; object=ch06-layer-cake; number=6.24; region=middle
\begin{proposition}\label{proposition:ch06-layer-cake}
若 $0<p<\infty$，则
\begin{align*}
\int_X|f|^p\,\mathrm d\mu=p\int_0^\infty\alpha^{p-1}\lambda_f(\alpha)\,\mathrm d\alpha.
\end{align*}
\end{proposition}
% SOURCE p=213; printed=198; object=ch06-layer-cake-proof; region=middle
\begin{proof}
若对某个 $\alpha>0$ 有 $\lambda_f(\alpha)=\infty$，则两边积分都为无穷。否则，先设 $f$ 为简单函数，此时 $\lambda_f$ 在 $\alpha\to0$ 时有界，并在 $\alpha$ 足够大时为零，故上面所说的分部积分可用。（这时也很容易直接验证公式。）对一般情形，取简单函数 $g_n\uparrow|f|$。所求结果对每个 $g_n$ 成立，再用\subcref{proposition:ch06-distribution-properties}{item:ch06-distribution-monotone-limit}和单调收敛定理，就得到对 $f$ 的结论。
\end{proof}
% SUPPLEMENT: id=v14-ch06-layer-cake-completion; parent=ch06-layer-cake; kind=proof-detail+motivation+hypotheses; origin=independent-derivation
\begin{remark}
\textbf{不依赖边界项的补证与条件说明（重构新增）。}若 $\lambda_f(a)=\infty$，左边至少为 $a^p\lambda_f(a)=\infty$；右边在 $(a/2,a)$ 上的被积函数已为无穷，所以也为无穷。若所有正水平都有限，非负简单函数 $s=\sum_{j=1}^m c_j\chi_{E_j}$ 的正值 $c_j>0$ 所对应的不交集合 $E_j$ 都有限测度。于是直接计算
\begin{align*}
p\int_0^\infty\alpha^{p-1}\lambda_s(\alpha)\,\mathrm d\alpha
&=\sum_{j=1}^m\mu(E_j)\,p\int_0^{c_j}\alpha^{p-1}\,\mathrm d\alpha\\
&=\sum_{j=1}^m c_j^p\mu(E_j)=\int s^p\,\mathrm d\mu.
\end{align*}
对 $0\leqslant g_n\uparrow|f|$，有 $g_n^p\uparrow|f|^p$ 和 $\lambda_{g_n}(\alpha)\uparrow\lambda_f(\alpha)$；在两边各使用单调收敛，即得完整公式。这条证明适用于任意测度 $\mu$，没有对不一定 $\sigma$ 有限的乘积空间擅自调用 Tonelli。

公式把每个点的高度 $|f(x)|^p$ 分解为从零到该高度的薄层贡献。$p>0$ 使 $p\int_0^c\alpha^{p-1}\,\mathrm d\alpha=c^p$ 并在零端可积，$p<\infty$ 保证是有限幂；可测性保证水平集定义。无须预设 $f\in L^p$，两边均可为无穷，也无需事先知道边界项消失；若 $f\in L^p$，边界项的确趋零，见\cref{exercise:ch06-distribution-end-limits}。
\end{remark}
% END SUPPLEMENT: id=v14-ch06-layer-cake-completion
% SOURCE p=213; printed=198; object=ch06-weak-lp-definition; region=lower
一种经常出现的 $L^p$ 空间变体如下。
\begin{definition}[弱 $L^p$]\label{definition:ch06-weak-lp}
若 $f$ 为 $X$ 上的可测函数，$0<p<\infty$，定义
\begin{align*}
[f]_p=\left(\sup_{\alpha>0}\alpha^p\lambda_f(\alpha)\right)^{1/p}.
\end{align*}
\textbf{弱 $L^p$} 是所有满足 $[f]_p<\infty$ 的 $f$ 所成的集合。
\end{definition}
% SOURCE p=213; printed=198; object=ch06-weak-lp-properties; region=lower
$[\cdot]_p$ 不是范数：容易验证 $[cf]_p=|c|[f]_p$，但三角不等式不成立。不过，弱 $L^p$ 是拓扑向量空间，见\cref{exercise:ch06-weak-lp-topology}。$L^p$ 与弱 $L^p$ 的关系如下：一方面，
\begin{align*}
L^p\subset\text{弱 }L^p,\qquad [f]_p\leqslant\|f\|_p.
\end{align*}
这只是 Chebyshev 不等式的另一种写法。另一方面，若在等于 $\|f\|_p^p$ 的积分 $p\int_0^\infty\alpha^{p-1}\lambda_f(\alpha)\,\mathrm d\alpha$ 中，以 $([f]_p/\alpha)^p$ 代替 $\lambda_f(\alpha)$，便得到常数倍的 $\int_0^\infty\alpha^{-1}\,\mathrm d\alpha$，它在零端与无穷端都发散，但恰好处于临界情形。只要在这两端对 $\lambda_f$ 的估计稍强一些，就能得到 $f\in L^p$。（另见\cref{exercise:ch06-weak-strong-inclusions}。）
% SUPPLEMENT: id=v14-ch06-weak-lp-motivation; parent=ch06-weak-lp-definition; kind=proof-detail+motivation+hypotheses; origin=independent-derivation
\begin{remark}
\textbf{定义的动机、识别约定与范数反例（重构新增）。}弱 $L^p$ 只要求所有水平集满足同一个幂衰减界 $\lambda_f(\alpha)\leqslant A^p\alpha^{-p}$，而不要求这些上界按层积分后仍有限。$[f]_p$ 正是这种 $A$ 的最小非负值。它与几乎处处代表无关；沿用本章 $L^p$ 的约定按几乎处处相等识别函数。若 $[f]_p=0$，每个 $\{|f|>1/n\}$ 都是零集，可数并给 $f=0$ 几乎处处，反向显然。可测性保证层集有测度，$0<p<\infty$ 决定正幂尺度；并不要求 $p\geqslant1$，因为这里还没有声称它是范数。对非零 $c$，$\lambda_{cf}(\alpha)=\lambda_f(\alpha/|c|)$，换上确界变量即得齐次性；$c=0$ 直接成立。

三角不等式一般确实失败。取两个点组成的计数测度空间，令 $a=2^{-1/p}$，$f=(1,a)$、$g=(a,1)$。由两个水平的直接计算，
\begin{align*}
[f]_p=[g]_p=\max\{1,2^{1/p}a\}=1,
\qquad [f+g]_p=2^{1/p}(1+a)=2^{1/p}+1>2.
\end{align*}
严格水平集不一定在阈值处取得上确界，但从下方逼近阈值即可得到这些值。拓扑向量空间的完整构造见习题解答。Chebyshev 给 $\alpha^p\lambda_f(\alpha)\leqslant\|f\|_p^p$，取上确界便证明所述连续包含估计；反向不能从临界的 $\alpha^{-1}$ 积分推出。
\end{remark}
% END SUPPLEMENT: id=v14-ch06-weak-lp-motivation
% SOURCE p=214; printed=199; object=ch06-weak-not-strong-example; region=top
\begin{example}\label{example:ch06-weak-not-strong}
在配备 Lebesgue 测度的 $(0,\infty)$ 上，$f(x)=x^{-1/p}$ 是属于弱 $L^p$ 而不属于 $L^p$ 的标准例子。
\end{example}
% SUPPLEMENT: id=v14-ch06-weak-example-proof; parent=ch06-weak-not-strong-example; kind=proof-detail+motivation+hypotheses; origin=independent-derivation
\begin{proof}
\textbf{直接核验（重构新增）。}对 $\alpha>0$，$\{x:x^{-1/p}>\alpha\}=(0,\alpha^{-p})$，故 $\lambda_f(\alpha)=\alpha^{-p}$，$[f]_p=1$。但 $\int_0^\infty|f(x)|^p\,\mathrm dx=\int_0^\infty x^{-1}\,\mathrm dx=\infty$，所以 $f\notin L^p$。
\end{proof}
\begin{remark}
\textbf{动机与条件说明（重构新增）。}这个函数在每个尺度都恰好达到弱型幂衰减，而层积分在两个端点呈对数发散。正半轴及 Lebesgue 测度使水平集的测度就是区间长度；$p>0$ 保证不等式解为所写区间。截取到 $(0,1)$ 可保留零端的发散，截取到 $(1,\infty)$ 可保留无穷端的发散，它们也说明习题中两个改进条件各自所控制的端点。
\end{remark}
% END SUPPLEMENT: id=v14-ch06-weak-example-proof
% SOURCE p=214; printed=199; object=ch06-truncation-introduction; region=upper
常常需要把函数写成“小部分”与“大部分”之和。下面的分解方式可为它们的分布函数给出简单公式。
% SOURCE p=214; printed=199; object=ch06-distribution-truncation; number=6.25; region=upper-middle
\begin{proposition}\label{proposition:ch06-distribution-truncation}
若 $f$ 为可测函数，$A>0$，令 $E(A)=\{x:|f(x)|>A\}$，并设
\begin{align*}
h_A&=f\chi_{X\setminus E(A)}+A(\operatorname{sgn}f)\chi_{E(A)},\\
g_A&=f-h_A=(\operatorname{sgn}f)(|f|-A)\chi_{E(A)}.
\end{align*}
则
\begin{align*}
\lambda_{g_A}(\alpha)=\lambda_f(\alpha+A),\qquad
\lambda_{h_A}(\alpha)=\begin{cases}\lambda_f(\alpha),&\alpha<A,\\0,&\alpha\geqslant A.\end{cases}
\end{align*}
\end{proposition}
% SOURCE p=214; printed=199; object=ch06-truncation-proof-omission; region=middle
原书将证明留给读者，见\cref{exercise:ch06-truncation-distributions}。
% SUPPLEMENT: id=v14-ch06-truncation-proof; parent=ch06-distribution-truncation; kind=omitted-proof+motivation+hypotheses; origin=independent-derivation
\begin{proof}
\textbf{原书省略证明的补全（重构新增）。}约定 $\operatorname{sgn}0=0$。当 $|f|\leqslant A$ 时，$h_A=f,g_A=0$；当 $|f|>A$ 时，$h_A=A\operatorname{sgn}f$、$g_A=(|f|-A)\operatorname{sgn}f$，且此时相位模为一。因此处处有
\begin{align*}
|h_A|=\min(|f|,A),\qquad |g_A|=(|f|-A)_+.
\end{align*}
对每个 $\alpha>0$，$\{|g_A|>\alpha\}=\{|f|>A+\alpha\}$。若 $0<\alpha<A$，则 $\{|h_A|>\alpha\}=\{|f|>\alpha\}$；若 $\alpha\geqslant A$，前一集合为空。取测度就得到全部公式。
\end{proof}
\begin{remark}
\textbf{动机与逐条件说明（重构新增）。}小部分保留原相位但把模长截在 $A$，大部分留下超过 $A$ 的高度，所以前者的分布函数在 $A$ 处截断，后者的分布函数向左平移。$f$ 可测使截断、相位和水平集可测，$A>0$ 提供正的分割高度；严格水平集说明在 $\alpha=A$ 时小部分的分布值为零。复值情形需要相位，不能简单写成实数的 $\min(f,A)$；所有公式在点值层面成立，无须任何测度有限性或可积性假设。
\end{remark}
% END SUPPLEMENT: id=v14-ch06-truncation-proof
\subfile{section-04-exercises.tex}
\end{document}
